\small
\begin{thebibliography}{99}
\bibitem{dysthe} K. B. Dysthe, Note on a modification to the nonlinear Schr\"odinger equation for application to deep water waves, \emph{Proc. R. Soc. Lond. A} \textbf{369} (1979), 105--114. User-supplied scan; especially equations (2.6), (2.18)--(2.21), (3.2)--(3.7), (4.1)--(4.3). \url{https://doi.org/10.1098/rspa.1979.0154}.
\bibitem{shemer} L. Shemer and B. Dorfman, Experimental and numerical study of spatial and temporal evolution of nonlinear wave groups, \emph{Nonlinear Processes in Geophysics} \textbf{15} (2008), 931--942. Potential/surface sign discussion on pp. 932--933. \url{https://doi.org/10.5194/npg-15-931-2008}.
\bibitem{trulsen} K. Trulsen, Weakly nonlinear and stochastic properties of ocean wave fields: Application to an extreme wave event, in \emph{Waves in Geophysical Fluids} (2006), Sect. 4.3. \url{https://doi.org/10.1007/978-3-211-69356-8_2}.
\bibitem{fedele} F. Fedele and D. Dutykh, Hamiltonian description and traveling waves of the spatial Dysthe equations, preprint (2011, revised 2012). Used here for the spatial equation and numerical coherent-wave context, not to resolve the temporal sign convention. \url{https://arxiv.org/abs/1110.3605}.
\bibitem{grande} R. Grande, K. M. Kurianski and G. Staffilani, On the nonlinear Dysthe equation, \emph{Nonlinear Analysis} \textbf{207} (2021), 112292. \url{https://arxiv.org/abs/2006.13392}.
\bibitem{mosincat} R. Mosincat, D. Pilod and J.-C. Saut, Global well-posedness and scattering for the Dysthe equation in $L^2(\mathbb R^2)$, \emph{J. Math. Pures Appl.} \textbf{149} (2021), 73--97. Formulation and small-data assumptions must be checked before application. \url{https://arxiv.org/abs/2007.01613}.
\bibitem{heller} G. Heller and C. Shao, Strichartz and multi-linear estimates for the one-dimensional periodic Dysthe equation, preprint (2021). \url{https://arxiv.org/abs/2112.12734}.
\bibitem{raissi} M. Raissi, P. Perdikaris and G. E. Karniadakis, Physics-informed neural networks: A deep learning framework for solving forward and inverse problems involving nonlinear partial differential equations, \emph{J. Comput. Phys.} \textbf{378} (2019), 686--707. \url{https://doi.org/10.1016/j.jcp.2018.10.045}.
\bibitem{causal} S. Wang, S. Sankaran and P. Perdikaris, Respecting causality for training physics-informed neural networks, \emph{Comput. Methods Appl. Mech. Engrg.} \textbf{421} (2024), 116813; preprint 2022. \url{https://arxiv.org/abs/2203.07404}.
\bibitem{failure} A. Krishnapriyan et al., Characterizing possible failure modes in physics-informed neural networks, \emph{NeurIPS} (2021). \url{https://arxiv.org/abs/2109.01050}.
\bibitem{fno} Z. Li et al., Fourier neural operator for parametric partial differential equations, \emph{ICLR} (2021); preprint 2020. \url{https://arxiv.org/abs/2010.08895}.
\bibitem{schober} C. M. Schober and A. Islas, Higher-Order Extensions of Weakly Viscous Dysthe Theory and a Phase-Lag Model for Nonlinear Mean-Flow Damping, user-supplied preprint, version 1, 5 August 2026. \url{https://arxiv.org/abs/2608.04473}.
\end{thebibliography}
\paragraph{How the literature is used.}
Classical PINNs and FNOs motivate the numerical architectures; they do not validate these methods for this exact water-wave equation. Causal and sequence-based PINN work \cite{causal,failure} motivates controlling temporal training difficulty. Our hard slab handoff is not identical to causal loss weighting. No general neural-method speed or accuracy claim is transferred from another PDE. The strongest Dysthe-specific evidence in this report is the explicitly labelled computation and source audit.
